\section{Dyskusja wyników}
\label{sec:dyskusja-wynikow}

Dotychczasowe eksperymenty pokazują, że na jakość strategii
w Ultimate Texas Hold'em wpływa nie tylko wybór algorytmu uczenia,
ale również sposób przedstawienia stanu i właściwości procesu
optymalizacji.

Najbardziej jednoznaczny efekt zaobserwowałem dla reprezentacji stanu
w DQN. Rozszerzenie RAW o cechy pokerowe w reprezentacji FEATURES
poprawiło wynik dla obu funkcji nagrody i efekt ten wystąpił
konsekwentnie dla pięciu sparowanych seedów treningowych. Wynik ten
wskazuje, że w badanym budżecie treningowym jawne dostarczenie
informacji o strukturze sytuacji pokerowej ułatwia DQN wykorzystanie
dostępnej obserwacji.

Dla REINFORCE analogicznego, stabilnego efektu reprezentacji stanu
nie zaobserwowałem. Najwyższe średnie EV uzyskało połączenie FEATURES
ze skalowaną nagrodą, jednak było ono jednocześnie konfiguracją
o największej zmienności pomiędzy seedami. Wynik średni nie
odzwierciedla więc w tym przypadku typowego rezultatu pojedynczego
treningu równie dobrze jak dla DQN.

Wpływ skalowania nagrody również nie był jednoznaczny. Dla DQN
korzystającego z reprezentacji RAW nagroda skalowana prowadziła do
wyższego EV, natomiast po zastosowaniu FEATURES różnica pomiędzy
funkcjami nagrody była niewielka. Diagnostyka norm gradientu
pokazała, że zmiana skali nagrody wyraźnie wpływała na obserwowaną
wielkość sygnału optymalizacyjnego, ale samo ograniczenie tej skali nie
gwarantuje poprawy końcowej polityki.

Porównując oba algorytmy przy podstawowym budżecie 50\,000 epizodów,
REINFORCE osiągnął wyższe średnie EV od DQN dla każdego
odpowiadającego sobie połączenia reprezentacji stanu i funkcji
nagrody. Nie oznacza to jednak jednoznacznej przewagi REINFORCE.
Wyniki tego algorytmu były bardziej zależne od seedu treningowego,
a większość przebiegów, osiemnaście z dwudziestu, prowadziła do
degeneracji polityki do jednej akcji niezależnej od stanu gry.

DQN osiągał niższe średnie EV, lecz proces treningowy był bardziej
powtarzalny. Również eksperyment z większym budżetem treningowym
wskazał na różnicę pomiędzy algorytmami. Wszystkie pięć modeli DQN
poprawiło wynik po zwiększeniu budżetu do 100\,000 epizodów,
natomiast dla REINFORCE poprawa była silnie zależna od konkretnego
seedu.

Żadna z badanych konfiguracji RL nie osiągnęła jakości strategii
RuleBased. Agent regułowy uzyskał w końcowej ewaluacji EV równe
$-0{,}0106$, podczas gdy najlepsze średnie wyniki podstawowych
konfiguracji RL pozostawały wyraźnie ujemne. Pokazuje to, że przy
zastosowanych architekturach i budżecie treningowym
prosta strategia wykorzystująca wiedzę dziedzinową pozostaje
znacznie skuteczniejsza.

Niemniej jednak nie oznacza to, że modele RL nie nauczyły się użytecznej
struktury problemu. Większość konfiguracji osiągnęła wynik wyższy od
agenta losowego, a analiza decyzji pokazała zależności pomiędzy
obserwowanym stanem i podejmowanymi akcjami. Wyniki sugerują więc,
że badane metody były zdolne do uczenia się strategii, lecz jakość
tej strategii pozostawała ograniczona.

REINFORCE
pozostał metodą bez funkcji bazowej, krytyki i dodatkowej regularyzacji
entropii. Reprezentacja FEATURES zawiera natomiast ręcznie
zaprojektowane cechy pokerowe, dlatego jej przewaga nad RAW pokazuje
wartość tej konkretnej informacji dziedzinowej, a nie uniwersalną
przewagę każdej reprezentacji o większej liczbie cech.

Z tego względu wyników nie traktuję jako rozstrzygnięcia o ogólnej
wyższości jednego algorytmu nad drugim. Pokazują one natomiast
wyraźnie, że w badanym środowisku DQN i REINFORCE reagują inaczej na
reprezentację stanu, skalę nagrody i zwiększenie budżetu treningowego.
Różnica pomiędzy średnią jakością polityki a stabilnością jej
uzyskiwania stanowi dość istotny rezultat przeprowadzonych eksperymentów.
